\scanpage{13}\mspage{10}

Or si $i\le n$, on aura
\[
 \left\{\ \Lambda\pi_{i0}=0,\quad \Lambda\pi_{i\alpha}=L^{\alpha-1}\unsure{W_{i-2\alpha}|w_{i-2\alpha}} L^{n-i+\alpha}\pi_{i\alpha}\quad\text{si }\alpha\ge1,\right.
 \tag*{(6.12)}
\]
\deleted{et si $j>n$, on aura}
\[
 \deleted{\Lambda\pi_{j,\alpha}=L}\tag*{\deleted{(6.13)}}
\]
\[
 \Lambda\pi_{2n-i,\alpha}=L^{n-i+\alpha-1}w_{i-2\alpha}\deleted{\unsure{\pi_{2n-i,\alpha}}}L^\alpha\pi_{2n-i,\alpha}\quad(\text{si }i\le n-1)
 \tag*{(6.13)}
\]
Les conditions envisagées impliquent évidemment (iv), : \unsure{prouvons|prouver} que (iv) implique (ii). Il suffit de prouver en fait que l'on a alors $\pi_i=\pi'_i$. \uline{Démonstration à fournir}.

\uline{Remarque} Voici un système de conditions déterminant de façon unique les éléments $\Lambda\in\unsure{E|\mathcal E}^{(2)}$, $\pi_{i,\alpha}\in E^{(0)}$ : en termes de $L$, et des $\pi_i^{\deleted{\unsure{*}}}$ (dont nous avons déjà une caractérisation intrinsèque par ailleurs, en termes de l'anneau gradué $E$) :

\margin{\deleted{\unsure{$\forall i$}}}
\[
 \left\{\begin{aligned}
 \text{\textcircled{\scriptsize 1}}\quad&(\Lambda^{n-i}L^{n-i}-1)\pi_i=0\qquad
 \text{\textcircled{\scriptsize 2}}\ \deleted{\unsure{\pi_{2n-i}}}(L^{n-i}\Lambda^{n-i}-1)\deleted{\unsure{\pi_i}}\pi_{2n-i}=0\\
 \text{\textcircled{\scriptsize 3}}\quad&\pi_{i,\alpha}=L^\alpha\deleted{\unsure{w_{i-2\alpha}}}\Lambda^{n-i+2\alpha}L^{n-i+\alpha}(1-\sum_{\alpha'>\alpha}\pi_{i,\alpha'})\\
 \text{\textcircled{\scriptsize 4}}\quad&\pi_{2n-i,\alpha}=L^{n-i}\pi_{i,\alpha}\Lambda^{n-i}\\
 &\deleted{\Lambda=\sum_{j,\alpha}\Lambda\pi_{j,\alpha}}\\
 \text{\textcircled{\scriptsize 5}}\quad&\Lambda\pi_{i,0}=0\ \unsure{\alpha}\\
 \text{\textcircled{\scriptsize 6}}\quad&\Lambda\pi_{i\alpha}=L^{\alpha-1}\Lambda^{n-i+2\alpha}L^{n-i+\alpha}\pi_{i\alpha}\quad\text{si }\alpha\ge1\\
 \text{\textcircled{\scriptsize 7}}\quad&\Lambda\pi_{2n-i,\alpha}=L^{n-i+\alpha-1}\Lambda^{n-i+2\alpha}L^\alpha\pi_{2n-i,\alpha}\quad\text{si }\unsure{\alpha|i}\le n-1\\
 \text{\textcircled{\scriptsize 8}}\quad&\Lambda=\sum_{j,\alpha}\Lambda\pi_{j,\alpha}
 \end{aligned}\right.
 \tag*{(6.14)}
\]

(dans ttes ces formules, on suppose $i\le n$).

Dans ces huit formules, on peut d'ailleurs remplacer les quatre dernières par les suivantes \inserted{plus simples} :
% archive sheet number: 12
